% ============================================================================
% CHAPTER 3 - METHODS / PARAMETERS
% File: docs/thesis/methods/group1_technical_potential.tex
% Group 1: technical potential (Stage F5). OQ-004, OQ-005, OQ-006, OQ-023.
% Theory: docs/thesis/theory/ch2_renewable_potential_foundations.tex
% Values must match config/parameters.json. References marked % verify are unchecked.
% ============================================================================

\section{Technical-Potential Parameters (Stage F5)}
\label{sec:params_group1}

\subsection{Scope and Cross-Cutting Decisions}
\label{subsec:g1_scope}

This section fixes the parameters that convert eligible land (Stage F3)
and the resource fields (Stages F2b and F4) into installable capacity and
capacity factor (Stage F5): land-utilization factor and power density
(OQ-004), wind turbine representation and losses (OQ-005), the Global
Wind Atlas inputs (OQ-006) and the PV temperature coefficient (OQ-023).
Costs are treated in Section~\ref{sec:params_group2}. Four decisions
apply to every parameter below.

Technology vintage. All technology parameters represent plants
commissioned in $t_{\mathrm{ref}}=2024$, the cohort of the cost data
(Section~\ref{sec:params_group2}). The same plant is evaluated
under the reference climate and under every projected climate member,
so that differences between members are attributable to climate. The
plant is therefore a fixed-vintage counterfactual
(Section~\ref{subsec:technology_vintage}).

Residual LUF. The LUF represents only constraints not encoded in the
Stage F3 exclusions (Section~\ref{subsec:luf_theory}). A justification
that refers to a constraint already mapped in F3 is not admissible. In
F3, cropland is excluded for PV in all three land scenarios and allowed
for wind; population density, protected areas and slope are excluded for
both.

Coupling to the land scenarios. LUF and power density are not sampled in
Stage F6. Their minimum, central and maximum values are paired with the
restrictive, central and permissive land-eligibility scenarios of
Stage F3. The three land configurations bound the joint effect of mapped
exclusions, residual utilization and density, without identifying the
contribution of each factor. Because LUF and PD scale all cells of a
country and technology uniformly, they affect quantities in GW and
TWh/yr (H5 and the comparison with national targets) but not the LCOE
of a cell or the regret ranking.

Electrical basis. PV capacity, PV power density, PV yield and PV unit
costs are in DC terms (MW$_{\mathrm{dc}}$, kWh/kW$_{\mathrm{dc}}$,
USD/kW$_{\mathrm{dc}}$). Wind capacity is the turbine nameplate rating.

\subsection{Land-Utilization Factor (OQ-004)}
\label{subsec:g1_luf}

Definition: dimensionless share of the Stage F3 eligible area expected
to host projects after residual, unmapped constraints. Enters:
Equation~\ref{eq:theory_capacity}. Scenario-coupled, not sampled.

The literature distinguishes total from available land and documents how
strongly assumed availability shifts potential estimates
\cite{lopez2012,ong2013,turkovska2021,turkovska2024}, but no source
reports a residual factor compatible with the F3 definition for Brazil,
India or Portugal. The values below are therefore author judgment
informed by that literature. Under U-07 this is Tier 3, and each value is
recorded in docs/LIMITATIONS.md.

\begin{table}[htbp]
    \centering
    \caption{Residual land-utilization factor (restrictive / central /
    permissive).}
    \label{tab:g1_luf}
    \small
    \begin{tabular}{lcc}
        \toprule
        Country & Solar PV & Onshore wind \\
        \midrule
        Brazil   & 0.10 / 0.25 / 0.35 & 0.05 / 0.15 / 0.30 \\
        Portugal & 0.10 / 0.25 / 0.35 & 0.05 / 0.15 / 0.30 \\
        India    & 0.05 / 0.10 / 0.25 & 0.05 / 0.075 / 0.15 \\
        \bottomrule
    \end{tabular}
\end{table}

The lower Indian values have different justifications for each
technology, because F3 treats cropland differently. For PV, cropland is
already excluded, so the Indian adjustment rests only on land-tenure
fragmentation and the difficulty of assembling large contiguous parcels
with clear title on the remaining land. For wind, cropland is eligible
in F3, and the adjustment also represents competition with agricultural
use on that land. Population density is not part of either
justification, since it is an F3 exclusion.

Source: author judgment informed by \cite{lopez2012,ong2013,turkovska2024}.
Tier 3 (value and range). Proxy: no. Price year: not applicable.

\subsection{Power Density (OQ-004)}
\label{subsec:g1_pd}

\subsubsection{Utility-Scale PV}

Definition: DC capacity per unit of fenced plant area. Unit:
MW$_{\mathrm{dc}}$/km$^2$. Enters: Equation~\ref{eq:theory_capacity}.
Scenario-coupled.

The PV baseline yield used by Stage F5 refers to fixed-tilt systems
(Section~\ref{subsec:g1_pv_thermal}), so only the fixed-tilt density is
used. Bolinger and Bolinger report a 2019 median of
0.87~MW$_{\mathrm{dc}}$/ha for fixed-tilt plants on the direct array
area, from more than 90\% of United States utility-scale plants, and
found plants checked in Chile, Australia and Europe within the same range
\cite{bolinger2022}. On total project area and an older fleet, Ong et
al.\ report lower densities \cite{ong2013}. The adopted values, 45, 50
and 75~MW$_{\mathrm{dc}}$/km$^2$, lie between these references; the
central value implies an array-to-site ratio of 0.57 relative to the
2019 array median of 87~MW$_{\mathrm{dc}}$/km$^2$. Modules installed in 2024
are more powerful than those of 2019, so the values are conservative. The tracking densities (30, 40 and
50~MW$_{\mathrm{dc}}$/km$^2$) are documented but not used, since no
tracking yield is modeled.
% verify: that the PV baseline yield layer is fixed-tilt at optimum angle (Global Solar Atlas PVOUT definition)

Source: \cite{bolinger2022,ong2013}, transferred from the United States
with the international check of \cite{bolinger2022} as rationale. Tier 2
(value and range). Proxy: no.

\subsubsection{Onshore Wind}

Definition: turbine nameplate capacity per unit of array spacing
envelope. Unit: MW/km$^2$. Enters: Equation~\ref{eq:theory_capacity}.
Scenario-coupled.

Density is geometric (Equation~\ref{eq:theory_wind_pd}) with the
reference turbine of Section~\ref{subsec:g1_turbine}
($P_{\mathrm{rated}}=3.37$~MW, $D=130$~m) and three spacing scenarios:
$6D\times10D$ (restrictive), $5D\times8D$ (central) and $4D\times7D$
(permissive), giving 3.32, 4.99 and 7.12~MW/km$^2$. Spacing scenarios
follow technical-potential and wind land-use studies
\cite{lopez2012,harrisonatlas2022,turkovska2024}. The density is the
same in the three countries because the turbine is the same.
% verify: that the cited studies use these spacing multiples

Empirical footprints of operating farms (Voronoi, convex hull) are used
only as a cross-check of the geometric values, since they depend on the
delineation method \cite{harrisonatlas2022,turkovska2021,turkovska2024}.

Source: geometric derivation from \cite{bortolotti2019} with spacing from
\cite{lopez2012,harrisonatlas2022}. Tier 2 (value and range). Proxy: no.

\subsection{Global Wind Atlas Inputs (OQ-006)}
\label{subsec:g1_gwa}

Stage F1 uses Global Wind Atlas version 3 (ERA5 2008--2017, WRF at about
3~km and WAsP to about 250~m) and extracts mean wind speed, Weibull $A$
and $k$ and air density at 100, 150 and 200~m
\cite{dtuwind2019gwa3release,gwaV3dataset}. The density layer derives
from CFSR interpolated to the atlas elevation, not from ERA5
\cite{gwaV3dataset}. The atlas capacity-factor layers are gross, computed
at 100~m for three Vestas 3.45~MW turbines (rotors 112, 126 and 136~m,
specific power about 350, 277 and 237~W/m$^2$) representing classes I,
II and III \cite{davis2023globalwindatlas}. Stage F1 does not download
them, and Stage F5 computes capacity factor independently. If the class
III layer is added as a check layer, it is compared with the gross
Stage F5 capacity factor at 100~m in the reference climate, since its
specific power is close to that of the GeoFREA turbine.
% verify: exact GWA release recorded in the F1 manifest (3.x)

Tier 1. OQ-006 is closed.

Data-validity bounds (OQ-053). At every extracted height, cells with
$A\le 0$, $A>25$~m/s, $k_{\mathrm{W}}\le 0$ or $k_{\mathrm{W}}>10$ are
treated as corrupted data and stop the run. These bounds test validity,
not plausibility: low shape parameters ($k_{\mathrm{W}}<1$, counted in
Stage F3) and high ones ($k_{\mathrm{W}}>4$, typical of steady trade-wind
regimes) are physical and are reported as distribution tails, not as
warnings. The numerical caps are author judgment (Tier 3).

\subsection{Wind Turbine Representation (OQ-005)}
\label{subsec:g1_turbine}

\subsubsection{Reference Turbine and Power Curve}

A single open reference turbine is used in every cell, country and
climate member: the IEA Wind Task 37 land-based reference turbine
(3.37~MW, rotor 130~m, hub 110~m, class IIIA, specific power about
254~W/m$^2$) \cite{bortolotti2019}. Its power curve is the tabulated
steady-state output of the published model (CCBlade), distributed under
the Apache-2.0 licence and archived with its version. The curve starts at
the cut-in speed of 3~m/s with a non-zero output; output is zero below
cut-in and above cut-out, and values are interpolated linearly between
tabulated points within that interval, with no extrapolation. The curve
is taken to refer to the air density of the model inputs.
% verify: reference air density in the turbine model inputs (expected 1.225 kg/m3)

The turbine is rated below the machines installed in 2024, but its
specific power is in the range of current low-wind rotors, and specific
power, not rating, sets the shape of the power curve and the density
(Equation~\ref{eq:theory_wind_pd}). A single curve keeps the turbine
identical across cells, countries and members. The per-cell IEC class
rule therefore does not select a curve; it is kept as a screening flag
(Section~\ref{subsec:g1_class}). This method change is recorded in the
METHODOLOGY changelog.

Tier 2. Proxy: no.

\subsubsection{Hub Height}

The Weibull parameters and air density are read from the 100~m layer in
the three countries, the layer nearest to the 110~m reference hub. The
10~m difference is not corrected with a shear exponent. The 150~m layer
is run once per country, in the reference climate, as a sensitivity case
for the taller towers of recent installations.

Tier 2. Proxy: no.

\subsubsection{Design-Class Screening Flag}
\label{subsec:g1_class}

For each cell, the reference mean speed at 100~m
(Equation~\ref{eq:wind_weibull_mean}) is compared with the class III
limit $V_{\mathrm{ave}}=7.5$~m/s (Table~\ref{tab:wind_iec_classes}).
Cells above it are flagged as exceeding the design envelope of the
reference turbine. The flag is computed once from the reference climate,
does not change across members and does not alter capacity factor or
eligibility. It reports the share of potential located where a real
project would use a stronger turbine. Extreme wind, turbulence and
tropical-cyclone exposure are not assessed \cite{larsen2022}.

Tier 1 (standard \cite{iec6140012025}). Proxy: no.

\subsubsection{Air-Density Correction}

The reference curve is stated at $\rho_0=1.225$~kg/m$^3$. The reference
density of a $0.01^\circ$ pixel $p$ is the GWA layer at 100~m, $\rho^{\mathrm{ref}}_p$ (the unit of
evaluation is the pixel, Section~\ref{subsec:g1_pixel_aggregation}).
For climate member $c$, density changes with temperature at constant
pressure:
\begin{equation}
    \rho_{i,c} = \rho^{\mathrm{ref}}_i\,
      \frac{T^{*}}{T^{*}+\Delta T_{i,c}},
    \qquad T^{*}=288.15~\mathrm{K},
    \label{eq:g1_density_member}
\end{equation}
where $\Delta T_{i,c}$ is the Stage F4 temperature change ($\Delta T=0$
in the reference). Using the standard-atmosphere temperature $T^{*}$ in
place of the local annual mean changes the density ratio by about 0.1\%
for local means between 0 and 30\,$^\circ$C and warming up to 6~K, which
is negligible against the other terms. The power curve is evaluated with
Equation~\ref{eq:wind_density_correction}; the output is not rescaled
again. One annual mean density per cell ignores the covariance between
density and wind speed, and the GWA density and wind fields come from
different reanalyses; both are limitations.

The correction is applied at every elevation; low densities on high
plateaus are physical and are not truncated. As a consistency check,
$\rho^{\mathrm{ref}}_i$ is compared with the International Standard
Atmosphere density at the cell elevation plus hub height, and cells
whose ratio falls outside 0.90--1.10 are reported. Whether very high
terrain is suitable for large turbines is a siting question (OQ-058):
until it is closed, wind cells above 3,000~m are excluded in Stage F3 in
all land scenarios, and their potential is reported as a sensitivity
case. Only India has eligible terrain above that elevation.

Tier 1 (ideal-gas law and IEC 61400-12-1 convention
\cite{iec61400122022}). Proxy: no.

\subsubsection{Aggregation from Pixel to Cell}
\label{subsec:g1_pixel_aggregation}

The Global Wind Atlas layers, the eligible fraction and the Stage F5 model
live on the $0.01^\circ$ pixels of Stage F2a, five by five in a
$0.05^\circ$ analysis cell. The power curve is not linear in the
Weibull parameters, so the capacity factor of the mean wind is not the mean
of the capacity factors. Stage F5 therefore evaluates the model of this
section (hub-height interpolation, density correction, Weibull integral,
retention) on every eligible pixel $p$ and aggregates to the cell $i$ with
the eligible area of the pixel as weight:
\begin{equation}
    \mathrm{CF}_{i,c} = \frac{\sum_{p\in i} w_p\,\mathrm{CF}_{p,c}}
      {\sum_{p\in i} w_p},
    \qquad w_p = f_p\,a_p,
    \label{eq:g1_wind_pixel_aggregation}
\end{equation}
where $f_p$ is the eligible fraction of the pixel in the land scenario and
$a_p$ its geodesic area. The weights add up to the eligible area of the
cell, so the energy of the cell,
$E_{i,c}=P_i\,\mathrm{CF}_{i,c}\,8760$, is the sum of the pixel energies at
the capacity density of the cell. $\Delta T_{i,c}$ and the wind-speed ratio
of Stage F4 are those of the cell and are constant over its pixels. Solar
is not affected: the model is linear in PVOUT and the two routes agree to
$10^{-16}$.

Before this definition, Stage F5 applied the curve to the eligible-area-weighted means of $A$, $k$ and
density of the cell. The difference was measured over every candidate cell of the
three countries in the reference climate
(Table~\ref{tab:g1_pixel_bias}). The cell value was below the pixel mean in
88\% to 91\% of the cells of India and Brazil, where wind speed varies
within a cell on rugged terrain, and symmetric in Portugal. The national energy moves by
less than 0.14\%, but the capacity factor of 6\% to 7\% of the
cells of India and Brazil was more than 1\% too low, in a few of them by 29\% to 53\%.
Because the cost of energy of a cell is inversely proportional to its energy, those cells would
have been ranked as costlier than they are.

\begin{table}[htbp]
    \centering
    \caption{Signed difference of the wind capacity factor computed on the cell means
    with respect to the eligible-area-weighted mean of the pixel capacity factors,
    $(\mathrm{CF}_{\mathrm{cell\,mean}}-\mathrm{CF}_{\mathrm{pixel}})/\mathrm{CF}_{\mathrm{pixel}}$,
    all candidate cells, reference member, central land scenario, before the change.}
    \label{tab:g1_pixel_bias}
    \small
    \begin{tabular}{lrrrrrrr}
        \toprule
        Country & Cells & Median & P5 & P95 & Beyond $\pm1\%$ & Largest & Energy \\
        \midrule
        Portugal & 3\,500   & $-0.003\%$ & $-0.24\%$ & $+0.26\%$ & 0.46\% & 2.0\%  & $+0.0006\%$ \\
        India    & 61\,464  & $-0.013\%$ & $-1.25\%$ & $+0.005\%$ & 5.8\% & 53.1\% & $-0.024\%$ \\
        Brazil   & 181\,306 & $-0.165\%$ & $-1.16\%$ & $+0.030\%$ & 6.8\% & 28.9\% & $-0.133\%$ \\
        \bottomrule
    \end{tabular}
\end{table}

Tier 1 (definition; exact aggregation). Proxy: no.

\subsubsection{Climate Forcing of the Resource}

The Stage F4 wind-speed factor modifies the reference distribution as
specified in METHODOLOGY (F4). Turbine, hub layer, density convention and
losses are identical across members. The wind ensemble uses the five
retained GCMs.

\subsubsection{Technical Losses}

Definition: retained fraction of gross energy after array wake,
unavailability, internal electrical losses and blade soiling and erosion,
$\eta$ in Equation~\ref{eq:wind_net_cf} (parameter eta\_loss in the
code, a retention factor). Curtailment is not modeled in GeoFREA;
long-term degradation is a Group 3 parameter. Enters: Stage F5 (central)
and Stage F6 (sampled, U-03).

\begin{table}[htbp]
    \centering
    \caption{Component losses of onshore wind and resulting retention.}
    \label{tab:g1_wind_losses}
    \small
    \begin{tabular}{lccc}
        \toprule
        Component loss & Minimum & Central & Maximum \\
        \midrule
        Array wake                & 3.0\% & 7.5\% & 12.0\% \\
        Availability              & 2.0\% & 3.0\% & 4.0\%  \\
        Electrical collection     & 1.0\% & 2.0\% & 3.0\%  \\
        Blade soiling and erosion & 1.0\% & 1.5\% & 2.0\%  \\
        \midrule
        Retention, $\eta$ & 0.932 & 0.866 & 0.803 \\
        \bottomrule
    \end{tabular}
\end{table}

The component intervals are author judgment organized by the taxonomy of
\cite{leefields2021}; no cited source reports these numbers for the three
countries. Stage F6 samples $\eta$ uniformly on $[0.803,\,0.932]$; Stage
F5 uses 0.866. The value is the same in the three countries.

Tier 3. Proxy: no.
% verify: compare component intervals with Lee and Fields (2021) before upgrading the tier

\subsection{PV Yield and Temperature Coefficient (OQ-023)}
\label{subsec:g1_pv_thermal}

\subsubsection{Yield Model}

Stage F5 takes the long-term specific yield of the Global Solar Atlas
(PVOUT, kWh/kW$_{\mathrm{dc}}$ per day) \cite{esmap2019gsa} as the reference, which already
includes thermal and system losses of a fixed-tilt system under the
reference climate. The capacity factor of cell $i$ in climate member $c$
follows Equation~\ref{eq:pv_perturbation}:
\begin{equation}
    \mathrm{CF}_{i,c} = \frac{\mathrm{PVOUT}_i}{24}\,
      \delta_{G,i,c}\,(1+\gamma\,\Delta T_{i,c}),
    \label{eq:g1_pv_cf}
\end{equation}
where $\delta_{G,i,c}$ is the Stage F4 ratio of surface irradiance and
$\Delta T_{i,c}$ the Stage F4 warming, both annual means
($\delta_G=1$, $\Delta T=0$ in the reference). The model uses no wind
input, so all six GCMs enter the solar ensemble.
% verify: PVOUT layer definition (fixed-tilt, optimum angle) and its loss assumptions in esmap2019gsa

\subsubsection{Temperature Coefficient of Maximum Power}

Definition: $\gamma$ in Equations~\ref{eq:pv_thermal_factor} and
\ref{eq:g1_pv_cf}. Unit: K$^{-1}$. Enters: Stage F5 (central) and Stage
F6 (sampled, U-03).

For the 2024 vintage the module class is n-type TOPCon, with
$\gamma=-0.0030$~K$^{-1}$ and a uniform range $[-0.0032,\,-0.0028]$
\cite{bamisile2025,itrpv2025}. PERC ($-0.0035$) and heterojunction
($-0.0026$) are documented as alternative technology cases and are not
run. The value is the same in the three countries. Because
Equation~\ref{eq:g1_pv_cf} is linear in $\gamma$, Stage F6 samples it
exactly from the stored terms $\mathrm{PVOUT}_i\,\delta_{G,i,c}/24$ and
$\Delta T_{i,c}$. In the reference climate $\gamma$ has no effect, since
the baseline yield already includes the thermal loss.

Tier 2. Proxy: no.
% verify: that bamisile2025 / itrpv2025 support the TOPCon interval

\subsection{Plausibility References for Stage F5}
\label{subsec:g1_plausibility}

Modeled capacity factors are compared, never calibrated, against two
kinds of reference. The first is the IRENA weighted average for plants
commissioned in 2024: 17.4\% for utility-scale PV and 34\% for onshore
wind worldwide, with country values where reported
\cite{irena2025costs2024}. The second is the national fleet capacity
factor computed from generation and average installed capacity of a
defined year \cite{epeAnuario2025,ceaGeneration,dgegRenewables}.
Both mix vintages, sites and curtailment
(Section~\ref{subsec:cost_basis}); a cell-level distribution is expected
to straddle them, not reproduce them. High values reported for recent
Brazilian wind projects represent the best sites, not the fleet.

\subsection{Summary}
\label{subsec:g1_summary}

\begin{table}[htbp]
    \centering
    \caption{Group 1 parameters. Scenario values are restrictive /
    central / permissive.}
    \label{tab:g1_summary}
    \small
    \renewcommand{\arraystretch}{1.15}
    \begin{tabularx}{\textwidth}{@{}>{\raggedright\arraybackslash}p{3.4cm}
        >{\raggedright\arraybackslash}p{1.7cm}
        >{\raggedright\arraybackslash}X
        >{\centering\arraybackslash}p{1.2cm}
        >{\centering\arraybackslash}p{0.7cm}@{}}
        \toprule
        Parameter & Unit & Value & Use & Tier \\
        \midrule
        LUF, PV, BRA and PRT & -- & 0.10 / 0.25 / 0.35 & scen. & 3 \\
        LUF, PV, IND & -- & 0.05 / 0.10 / 0.25 & scen. & 3 \\
        LUF, wind, BRA and PRT & -- & 0.05 / 0.15 / 0.30 & scen. & 3 \\
        LUF, wind, IND & -- & 0.05 / 0.075 / 0.15 & scen. & 3 \\
        PD, PV fixed-tilt & MW$_{\mathrm{dc}}$/km$^2$ & 45 / 50 / 75 & scen. & 2 \\
        PD, wind & MW/km$^2$ & 3.32 / 4.99 / 7.12 & scen. & 2 \\
        Reference turbine & -- & IEA 3.4-130 (3.37 MW, 130 m) & fixed & 2 \\
        Resource layer & m & 100 (150 sensitivity) & fixed & 2 \\
        Class flag threshold & m/s & 7.5 & flag & 1 \\
        $\rho_0$; $\rho^{\mathrm{ref}}$; $T^{*}$ & kg/m$^3$; K & 1.225; GWA layer; 288.15 & fixed & 1 \\
        $\eta$ (retention) & -- & 0.866, U[0.803, 0.932] & F6 & 3 \\
        $\gamma$ (TOPCon) & K$^{-1}$ & $-0.0030$, U[$-0.0032$, $-0.0028$] & F6 & 2 \\
        \bottomrule
    \end{tabularx}
\end{table}

No parameter in this group is a proxy, and none carries a price year.

\subsection{Limitations}
\label{subsec:g1_limitations}

The residual LUF has no primary source in any country and dominates the
spread of potential in GW and TWh/yr; results for H5 and for the
comparison with national targets are conditional on it. PV density is
transferred from United States plants of 2019 and is likely
conservative for 2024 modules. The PV yield is a fixed-tilt baseline
perturbed with annual means, so changes in the timing of warming and the
effect of wind and irradiance on cell temperature are not represented,
and tracking systems are not modeled. Wind losses rest on judgment. A
single class III turbine is applied everywhere, rated below 2024
machines; the class flag quantifies the share of potential where a
stronger turbine would be chosen. The 100~m layer represents a 110~m
hub. Air density uses annual means and a reanalysis different from the
wind field. The climate factors of a wind cell are those of the GCM
cell and are constant over its pixels. Wind potential above 3,000~m, present only in India, is excluded
pending OQ-058. Inverter clipping and curtailment are not modeled. The plant
is a 2024-vintage counterfactual evaluated under mid-century climate.
